\begin{prob}
    Determine the worst case time complexity of each of the recursive algorithms below.
    In each case, state the recurrence relation describing the runtime. Solve the recurrence
    relation, either by unrolling it or showing that it is the same as a recurrence we have
    encountered in lecture.

    \begin{subprobset}
        \begin{subprob}
            \inputminted{python}{\thisdir/include/part_a.py}

            \begin{soln}
                Note that the slicing in \python{numbers[:middle]} takes linear time
                in the length of the array, so the recurrence relation is:
                \[
                    T(n) = \Theta(n) + 2T(n/2)
                \]
                This is the same recurrence relation as for mergesort, and the solution
                is $\Theta(n \log n)$.
            \end{soln}
        \end{subprob}


        \begin{subprob}

            \inputminted{python}{\thisdir/include/part_b.py}

            \begin{soln}
                No slicing is being done, so a constant amount of time is required outside
                of the recursive calls. Therefore the recurrence relation is:
                \[
                    T(n) = 3T(n/3) + \Theta(1)
                \]

                We have not seen this recurrence before, so we must solve it by unrolling.
                It isn't so different from $T(n) = 2 T(n/2) + 1$, which you may have seen
                in lecture or in discussion. Both have $\Theta(n)$ as their answer.
            \end{soln}
        \end{subprob}

        \begin{subprob}
            In this problem, remember that \python{//} performs \emph{flooring division},
            so the result is always an integer. For example, \python{1//2} is zero.
            \python{random.randint(a,b)} returns a random integer in $[a,b)$ in constant time.

            \inputminted{python}{\thisdir/include/part_c.py}

            \begin{soln}
                The work outside of the recursive calls takes linear time. There
                are three recursive calls, each on problems of size (roughly) $n/3$.
                Therefore the recurrence is:
                \[
                    T(n) = 3T(n/3) + \Theta(n)
                \]
                We can solve this recurrence by unrolling it. We'll find that the solution is
                $\Theta(n \log n)$. Note the similarity between this recurrence and that
                for mergesort, and that both have the same solution.
            \end{soln}
        \end{subprob}

    \end{subprobset}

\end{prob}
